%AdaPop is a gradient-based weight-editing method: this family directly modifies parametric memory at training time, scales to arbitrary knowledge types without architectural assumptions, and is the class in which the popularity gap-differential resistance of facts to equal gradient force-manifests most directly. Methods such as GD, NPO, and WGA require manual tuning of unlearning hyperparameters like the retain coefficient $\alpha$: too small and the model loses general utility; too large and target facts survive erasure. AdaPop replaces this manual search with two components: a popularity-dependent exponent $\beta_i$ that calibrates the per-fact unlearning signal from an external metric, and a dual-ascent controller that adjusts $\alpha$ automatically during training.
We frame targeted unlearning as a constrained optimisation: increase the forget-set loss while bounding the retain-set drift. The proposed AdaPop solves this by combining (i) per-fact popularity scaling via a power-law exponent computed from an external popularity proxy that reweights per-token ascent signals, and (ii) a dual-ascent retain controller that updates the retain penalty at epoch granularity. Together, these components close the ``popularity gap'' by (a) concentrating the ascent signal on the highest-confidence answer tokens for rare facts, where encoding is narrow and over-erasure risk is high, and (b) distributing the ascent signal uniformly across all answer tokens for popular facts, where broad parametric encoding demands wide gradient coverage.



\subsection{The Popularity Gap}

% TODO: It should be noted that log-probability does not accurately reflect the actual popularity of a fact, as it is subject to noise from prepositions and particles, which typically have a high probability.
Parametric memorisation scales with training-corpus frequency~\cite{hartmann2024undesirable, kandpal2023large}. A single shared $\alpha$ (as in GD) or a confidence-based weight (as in WGA) cannot simultaneously address both regimes: popular facts resist erasure while rare facts risk collateral damage from the same gradient magnitude. AdaPop assigns each fact an independent exponent $\beta_i$ derived from an external popularity score, so that the unlearning signal is calibrated to fact-level memorisation depth rather than to the model's current output distribution.



\subsection{Problem Formulation}

Let $\mathcal{F}=\{(x_i,y_i,s_i)\}_{i=1}^{N_F}$ be the \emph{forget} set, where each entry consists of a question-answer pair $(x_i, y_i)$ and a corresponding popularity score $s_i$ derived from a global metric such as Wikidata citelinks count (\texttt{pop\_sum}). Let $\mathcal{R}=\{(x_j,y_j)\}_{j=1}^{N_R}$ be the \emph{retain} set. We consider an autoregressive model $q_\theta$ defined by the conditional probability $q_\theta(y_t \mid x, y_{<t})$ of a token $y_t$ given its prefix. For a token position $t$ in sequence $i$, we define the negative log-likelihood (NLL)
\begin{equation}
    \ell_{i,t}(\theta) = -\log q_\theta(y_{i,t} \mid x_i, y_{i,<t})
\end{equation}
and the model's confidence
\begin{equation}
    p_{i,t}(\theta) = \exp(-\ell_{i,t}(\theta)).
\end{equation}

The sets $\Omega_{\mathcal{F}}$ and $\Omega_{\mathcal{R}}$ represent the indices of valid tokens not masked by padding or special tokens. The objective of AdaPop is to increase the loss on $\mathcal{F}$ while bounding the drift on $\mathcal{R}$. We frame this as a constrained optimisation problem (maximise forgetting subject to $\delta_k \le \varepsilon$) and enforce it via a dual-ascent controller (§\ref{sec:dual_ascent}). A fixed retain coefficient $\alpha$, as in GD~\cite{liu2022continual} and WGA~\cite{wang2025rethinking}, requires dataset-specific grid search; the constrained formulation eliminates this by adjusting $\alpha$ automatically in response to observed retain drift. Full design rationale is in Appendix~\ref{app:design_motivation}.



\subsection{Popularity-Aware Forget Objective}

Each sample $i$ in the forget set is assigned an exponent $\beta_i$ based on its popularity score $s_i$:
\begin{equation}
    \beta_i = a \cdot s_i^{-b},
\end{equation}
A power-law mapping is natural here because parametric memorisation itself scales as a power law with corpus frequency~\cite{kandpal2023large}: since $s_i$ proxies that frequency, $\beta_i \propto s_i^{-b}$ mirrors the underlying relationship. This choice of power law compresses the proxy’s wide dynamic range into a bounded exponent that places rare facts in a self-limiting regime ($\beta>1$) and popular facts in a pressure-sustaining regime ($\beta<1$), improving forget/retain tradeoffs. The coefficients $a$ and $b$ are proxy-specific (see Appendix~\ref{app:derivation} for derivation). $\beta(s)$ is monotonically decreasing: as popularity increases, $\beta_i$ decreases, distributing the ascent signal more broadly across tokens for popular facts. For rare facts, a larger $\beta_i$ concentrates the update on the highest-confidence tokens, limiting collateral drift. Values are clipped to $[0.05, 2.0]$ in all experiments to avoid numerical instability.

The per-token unlearning weight is:
\begin{equation}
w_{i,t} = \left(\mathrm{stopgrad}(p_{i,t}(\theta))\right)^{\beta_i}.
\end{equation}
The forget loss is:
\begin{equation}
L_f(\theta) = -\frac{1}{|\Omega_{\mathcal{F}}|} \sum_{(i,t) \in \Omega_{\mathcal{F}}} w_{i,t} \ell_{i,t}(\theta).
\end{equation}

%\subsection{Implementation Note}

The weighting $w_{i,t}$ is applied only during training, not at inference. It is computed only over answer tokens $y_i$; question tokens $x_i$ are masked and do not contribute to the loss, so the model's ability to parse the query remains unaffected.





\subsection{Dual-Ascent and Auto-Balancing}\label{sec:dual_ascent}

AdaPop treats $\alpha$ as a dual variable enforcing the constraint $\delta_k \le \varepsilon$, where $\delta_k$ is the one-sided relative drift of the retain loss $R_k$ at epoch $k$:
\begin{equation}
\delta_k = \max\left( 0, \frac{R_k - R_{\mathrm{ref}}}{\max(R_{\mathrm{ref}}, \epsilon_0)} \right),
\end{equation}
$R_{\mathrm{ref}}$ is the retain loss measured at the end of the first epoch, and $\epsilon_0 > 0$ is a small constant for numerical stability.

\paragraph{Epoch-Level Dual Update.}
Let $\alpha_k = \alpha_0 + \lambda_k$. The dual variable $\lambda$ is updated via projected gradient ascent once per epoch (not per step) to reduce oscillation due to noisy per-batch retain losses:
\begin{align}
    \lambda_{k+1} &= \Pi_{[0, \lambda_{\max}]}\left( \lambda_k + \eta_\lambda (\delta_k - \varepsilon) \right) \\
    \alpha_{k+1} &= \Pi_{[0,\, \alpha_0 + \lambda_{\max}]}(\alpha_0 + \lambda_{k+1}),
\end{align}
where the dual updates enforce the retain constraint on the popularity-weighted forget loss (rather than the unweighted or entropic formulations), aligning the feedback signal with the actual gradient mass distribution and enabling the controller to find operating points with deeper, targeted forgetting. We empirically found both choices important for stability (see Fig.~\ref{fig:ablation} in the Appendix).

\paragraph{Per-Step Popularity Modulation (optional).}
AdaPop supports an optional fine-grained control mode that scales the effective retain coefficient per batch. Since batches containing popular facts (small $\beta$) generate aggressive unlearning signals, the effective retain coefficient $\alpha_{\mathrm{eff}} = \alpha \cdot \text{scale}$ can be boosted to compensate:
\begin{equation}
\text{scale} = \mathrm{clip}\!\left(\frac{1}{1 + g \cdot \bar{\beta}},\; s_{\min},\; s_{\max}\right)
\end{equation}
where $g \geq 0$ is a gain coefficient, $\bar{\beta}$ is the batch-mean exponent, and $[s_{\min}, s_{\max}]$ clips the scale to a stable range. When $g = 0$, $\text{scale} = 1$, and the batch-level modulation is disabled. In all experiments reported in this paper, $g = 0$ (i.e., per-step modulation is off); the dual-ascent epoch-level controller alone provides the retain's balance.

\begin{algorithm}[t]
\caption{AdaPop}
\label{alg:adapop}
\begin{algorithmic}[1]
\REQUIRE Forget set $\mathcal{F}=\{(x_i,y_i,s_i)\}$ with popularity scores $s_i$, retain set $\mathcal{R}$, model $q_\theta$, initial retain weight $\alpha_0$, dual step size $\eta_\lambda$, tolerance $\varepsilon$, anchor points $(s_r, s_p)$
\ENSURE Updated parameters $\theta$

\STATE Compute popularity exponents:
\(
\beta_i = \mathrm{clip}(a\, s_i^{-b}, \beta_{\min}, \beta_{\max})
\)
where coefficients $a,b$ are derived from anchor constraints $\beta(s_r)$ and $\beta(s_p)$.

\STATE Initialize $\lambda \leftarrow 0$, $\alpha \leftarrow \alpha_0$

\FOR{epoch $k = 1$ to $E$}

    \FOR{batch $B$ sampled from $\mathcal{F} \cup \mathcal{R}$}

        \FOR{each forget example $i$ and answer token $t$ in $B$}
            \STATE $p_{i,t} \leftarrow \mathrm{stopgrad}(\exp(-\mathrm{NLL}_{i,t}(\theta)))$
            \STATE $w_{i,t} \leftarrow p_{i,t}^{\beta_i}$
        \ENDFOR

        \STATE Compute weighted forget loss:
        \[
        L_f = - \frac{1}{|B_F|} 
        \sum_{(i,t)\in B_F} w_{i,t}\,\mathrm{NLL}_{i,t}(\theta)
        \]

        \STATE Compute retain loss $L_r$ over retain tokens

        \STATE Total loss:
        \(
        L = L_f + \alpha L_r
        \)

        \STATE Update model parameters $\theta$ using gradient descent (LoRA updates)

    \ENDFOR

    \STATE Compute retain drift:
    \(
    \delta_k = \frac{\max(0, R_k - R_{\text{ref}})}{\max(R_{\text{ref}}, \epsilon_0)}
    \)

    \STATE Dual update:
    \(
    \lambda \leftarrow \mathrm{proj}_{[0,\lambda_{\max}]}\big(\lambda + \eta_\lambda(\delta_k - \varepsilon)\big)
    \)

    \STATE Update retain weight:
    \(
    \alpha \leftarrow \mathrm{proj}_{[0,\alpha_0+\lambda_{\max}]}(\alpha_0 + \lambda)
    \)

\ENDFOR

\RETURN $\theta$
\end{algorithmic}
\end{algorithm}

Two design choices in the proposed AdaPop (Alg.~\ref{alg:adapop}) are worth highlighting. The dual update runs at epoch granularity: per-batch retain losses are noisy, and epoch-level feedback reduces oscillation while still responding to retain drift within a few hundred gradient steps. The controller uses a popularity-weighted forget loss rather than an unweighted formulation, aligning the feedback signal with the actual gradient mass distribution and enabling the controller to find operating points with deeper, targeted forgetting. Ablations in Appendix~\ref{app:ablation} confirm both choices are necessary.



%% take from other/ada_pop.tex

% In this section, we detail the \textbf{AdaPop} (Adaptive Popularity) framework. Our approach moves beyond local model confidence by integrating a global popularity metric and an automated dual-ascent controller to maintain the balance between knowledge erasure and model utility.

% \subsection{Motivation and Setting}

% Let $\mathcal{F}=\{(x_i,y_i,s_i)\}_{i=1}^{N_F}$ be the \emph{forget} set, where each entry consists of a question-answer pair $(x_i, y_i)$ and a corresponding popularity score $s_i$ derived from a global metric such as \texttt{pop\_sum}. Let $\mathcal{R}=\{(x_j,y_j)\}_{j=1}^{N_R}$ be the \emph{retain} set used to preserve general model performance. We consider an autoregressive model $q_\theta$ where the conditional probability of a token $y_t$ given its prefix is $q_\theta(y_t \mid x, y_{<t})$. For any token position $t$ in sequence $i$, we define the negative log-likelihood (NLL) and the model's confidence as:
% \begin{align}
%     \ell_{i,t}(\theta) &= -\log q_\theta(y_{i,t} \mid x_i, y_{i,<t}) \\
%     p_{i,t}(\theta) &= \exp(-\ell_{i,t}(\theta))
% \end{align}
% The sets $\Omega_{\mathcal{F}}$ and $\Omega_{\mathcal{R}}$ represent the indices of valid tokens not masked by padding or special tokens. The objective of AdaPop is to increase the loss on $\mathcal{F}$ while minimizing the drift on $\mathcal{R}$ through a weighted gradient ascent objective where the retain coefficient $\alpha$ is dynamically adjusted.

% \subsection{Popularity-Aware Forget Objective}

% A core contribution of AdaPop is the use of a \textbf{popularity-dependent exponent} $\beta_i$ to scale the unlearning signal. Each sample $i$ in the forget set is assigned an exponent based on the popularity score $s_i$ of the corresponding fact. The exponent is defined as a general function of popularity:
% \[
% \beta_i = f(s_i).
% \]
% In practice, this function is often defined empirically to control the strength of the update. For example, one possible form used in our implementation is:
% \[
% \beta_i = 58.7 s_i^{-0.796},
% \]
% where the constant $58.7$ and the exponent $-0.796$ were chosen empirically to balance the influence across the popularity spectrum. To ensure training stability and prevent gradient explosion, $\beta_i$ is typically clipped to a stable interval (e.g., $[0.05, 2.0]$). Alternatively, a constant exponent $\beta_i = \beta_{\mathrm{const}}$ may be used to treat all facts with equal importance.



% The relationship between popularity and the final unlearning weight $w_{i,t}$ is governed by the detached token weight:
% \[
% w_{i,t} = \left(\mathrm{stopgrad}(p_{i,t}(\theta))\right)^{\beta_i}.
% \]
% This formulation creates a critical inverse relationship. Popular facts (high $s_i$) result in a smaller $\beta_i$, which amplifies and "flattens" the unlearning signal. This ensures that even tokens with moderate confidence contribute significantly to the gradient ascent, overcoming the strong "memory inertia" of deeply embedded knowledge. Conversely, rare facts (low $s_i$) receive a larger exponent, which "sharpens" the update to concentrate only on the most certain tokens, preventing unnecessary drift on niche information. The final forget loss is defined as:
% \[
% L_f(\theta) = -\frac{1}{|\Omega_{\mathcal{F}}|} \sum_{(i,t) \in \Omega_{\mathcal{F}}} w_{i,t} \ell_{i,t}(\theta).
% \]

% \subsection{Dual-Ascent and Fine-Grained Retain Control}

% To maintain stability without manual hyperparameter tuning, AdaPop employs a two-tier control system for the retain coefficient $\alpha$.

% \paragraph{Epoch-Level Dual Update.} 
% AdaPop treats $\alpha$ as a dual variable to satisfy the constraint $\delta_k \le \varepsilon$, where $\delta_k$ is the one-sided relative drift of the retain loss $R_k$ at epoch $k$:
% \[
% \delta_k = \max\left( 0, \frac{R_k - R_{\mathrm{ref}}}{\max(R_{\mathrm{ref}}, \epsilon_0)} \right).
% \]
% Let $\alpha_k = \alpha_0 + \lambda_k$. The dual variable $\lambda$ is updated via projected gradient ascent:
% \begin{align}
%     \lambda_{k+1} &= \Pi_{[0, \lambda_{\max}]}\left( \lambda_k + \eta_\lambda (\delta_k - \varepsilon) \right) \\
%     \alpha_{k+1} &= \Pi_{[0,\, \alpha_0 + \lambda_{\max}]}(\alpha_0 + \lambda_{k+1})
% \end{align}

% \paragraph{Per-Step Popularity Modulation.} 
% Beyond epoch-level updates, AdaPop employs fine-grained control to protect the model from sudden collapse. Since batches containing popular facts (small $\beta$) generate aggressive unlearning signals, we proactively scale the effective retain coefficient $\alpha_{\mathrm{eff}}$ for each batch:
% \[
% \text{scale} = \frac{1}{1 + \text{gain} \cdot \bar{\beta}}
% \]
% This scaling strengthens the "shield" of the retain loss in real-time when the model is performing the loud, wide updates required for popular facts, ensuring that the unlearning process does not cause immediate collateral damage to general knowledge.